\paragraph{\trFibrationWithSingularities{}。}\label{sub:fibrations-with-singularities}
我们将考察映射 $\pi\colon E\to B$，其中 $E$ 与 $B$ 是\trExactSymplectic\trCorneredBoundary\trManifold，$B$ \trConnected，并且 $\dim(E)\ge\dim(B)$，而 $\pi$ 是 $(I_E,I_B)$-\trHolomorphic{}的。我们再作三条附加的假设，它们合在一起保证这些映射沿边界有合理的表现。
\begin{itemize}
\item \emph{对 $\partial B$ 的\trTransversality。}设 $C\subset\partial B$ 是任何一个\trBoundaryStratum。在每一个使得 $y=\pi(x)\in C$ 的点 $x\in E$ 处，我们有 $D\pi(TE_x)+TC_y=TB_y$。
\end{itemize}

这显然蕴含 $\pi^{-1}(C)$ 是 $E$ 的一个与 $C$ 同\trCodimension{}的\trBoundaryStratum。这类\trStratum{}的并便是 $\partial^v E=\pi^{-1}(\partial B)$，我们称之为 $E$ 的\textbf{\trVertical{}边界}（vertical boundary）。所有不包含于 $\partial^v E$ 的\trBoundaryFace{}的并，称为 $E$ 的\textbf{\trHorizontal{}边界}（horizontal boundary），并记作 $\partial^h E$。
\begin{itemize}
\item \emph{沿 $\partial^h E$ 的\trRegularity。}若 $F$ 是 $E$ 的一个不包含于 $\partial^v E$ 的\trBoundaryFace，则 $\pi|F\colon F\to B$ 是一个\trSmooth\trFibration。
\end{itemize}

这蕴含任何\trFiber{}$E_y$ 在其边界 $\partial E_y=E_y\cap\partial^h E$ 附近都是\trSmooth{}的。到此为止，讨论一直集中于\trDifferential\trTopology，但我们还需要照顾辛的方面。在任一点 $x\in E$ 处，\trTangentSpace{}都\trSplitting{}成
\begin{equation}\label{eq:tangent-splitting}
TE_x=TE^h_x\oplus TE^v_x
\end{equation}
其中\trVertical{}部分是 $TE^v_x=\ker(D\pi_x)$，而\trHorizontal{}部分 $TE^h_x$ 是它关于 $\phi_E$ 的\trOrthogonal{}补（\trEquivalence{}地，对于由 $\phi_E$ 与 $I_E$ 得到的\trMetric{}亦然；这样的\trSplitting{}我们已经在\cref{sub:equivariant-transversality}中简要地遇到过，不过是在一个纯粹技术性的角色中）。我们的最后一条要求是
\begin{itemize}
\item \emph{$\partial^h E$ 的\trHorizontal{}性。}若 $x$ 落在某个不包含于 $\partial^v E$ 的\trBoundaryFace{}$F$ 中，则 $TE^h_x\subset TF_x$。
\end{itemize}

一个满足上述三条条件的映射 $\pi\colon E\to B$，称为一个\textbf{（\trExactSymplectic{}的）\trFibrationWithSingularities}（fibration with singularities）。

让我们考察几个简单的推论。记 $\mathrm{Crit}(\pi)\subset E$ 为\trCriticalPoint{}所成之集，而 $\mathrm{Critv}(\pi)=\pi(\mathrm{Crit}(\pi))\subset B$ 为\trCriticalValue{}所成之集。如上文已经提到的，沿 $\partial^h E$ 的\trRegularity{}意味着 $\mathrm{Crit}(\pi)\subset E\setminus\partial^h E$。在 $\dim(B)=2$ 的情形，对 $\partial B$ 的\trTransversality{}还蕴含 $\mathrm{Crit}(\pi)\cap\partial^v E=\varnothing$，但对更高维的底空间，这不再成立。
\begin{lemma}\label{thm:fibre-exact-symplectic}
若 $y\notin\mathrm{Critv}(\pi)$，则\trFiber{}$E_y$（装备着限制 $\phi_{E_y}=\phi_E|_{E_y}$、$\theta_{E_y}=\theta_E|_{E_y}$、$I_{E_y}=I_E|_{E_y}$）是一个\trExactSymplectic\trCorneredBoundary\trManifold。
\end{lemma}
\begin{proof}
唯一需要检验的非平凡之事是：负\trLiouville\trVectorField{}沿 $\partial E_y=E_y\cap\partial^h E$ \trIncoming。我们知道，$E$ 的负\trLiouville\trVectorField{}沿每一个\trBoundaryFace{}$F$ 都严格\trIncoming。现在假设 $F$ 不属于 $\partial^v E$。由于 $TE^h$ 平行于 $F$，我们这个\trVectorField{}的 $TE^v$ 分量沿 $F$ 也严格\trIncoming。而这个分量恰好就是 $E_y$ 的负\trLiouville\trVectorField，这就给出了所要的结果。
\end{proof}

由 \eqref{eq:tangent-splitting} 给出的\trSplitting{}在 $B\setminus\mathrm{Critv}(\pi)$ 之上定义了一个\trConnection。由于 $TE^h$ 平行于 $\partial^h E$，这个\trConnection{}具有\trWellDefined{}的\trParallelTransport{}映射，这意味着：对每条路径 $\beta\colon[0;1]\to B\setminus\mathrm{Critv}(\pi)$，我们都得到一个特异的\trDiffeomorphism（它\trSmooth{}地依赖于路径的选取）
\begin{equation}\label{eq:parallel-transport-map}
h_\beta\colon E_{\beta(0)}\longrightarrow E_{\beta(1)}.
\end{equation}
一个标准的计算表明，$h_\beta$ 实际上是一个\trExactSymplectic\trCorneredBoundary\trManifold{}之间的\trIsomorphism。取定 $\pi$ 的一个\trRegularValue{}$*$，把相应的\trFiber{}记为 $M=E_*$，并令 $B'\subset B\setminus\mathrm{Critv}(\pi)$ 是 $*$ 的一个小邻域。利用一组适当的\trParallelTransport{}映射，便得到 $B'$ 之上的一个局部\trTrivialization{}$E'=B'\times M\to E$。把 $\theta_E$、$\phi_E$ \trPullback{}所得的\trDifferential\trForm{}$\theta_{E'}$、$\phi_{E'}$ 呈如下形状：
\begin{equation}\label{eq:local-trivialization-forms}
\theta_{E'}=\theta_M+\kappa+d(\text{某个函数}),\qquad \phi_{E'}=\phi_M+d\kappa,
\end{equation}
其中 $\kappa$ 是一个在 $(TE')^v$-方向上消失的 1-\trForm。它可以被看作底空间上一个取值于函数的 1-\trForm，即 $\kappa\in\Omega^1(B',C^\infty(M,\RR))$。实际上，如果从函数过渡到\trHamiltonianVectorField，那么 $\kappa$ 恰好是 $(TE')^h$ 的\trConnection{} 1-\trForm。为看到这一点，取 $B'$ 上某个\trVectorField{}$Y$，并令 $Y^h=(Y,X)$ 为它到 $E'$ 的唯一\trHorizontal\trLift。在 Cartan 公式 $i_{Y^h}(d\kappa)=L_{Y^h}\kappa-d(i_{Y^h}\kappa)$ 中，\trLie{}导数又是一个在\trFiber{}方向上消失的 1-\trForm。因此，$Y^h$ 的定义性质可以写成如下形式：
\begin{equation*}
i_X\phi_M-d(i_Y\kappa)=0\qquad\text{在每根\trFiber{}}\ \{y\}\times M\ \text{上}.
\end{equation*}
于是，$X$ 逐\trFiber{}地是函数 $-\kappa(Y)$ 的\trHamiltonianVectorField，这就确立了 $\kappa$ 的上述性质。总结起来，这些考量表明：在 $B$ 的\trRegular{}部分之上，我们所拥有的是一个带优先\trConnection{}的 $\mathrm{Aut}(M)$-\trFiberBundle。注意，$B$ 上的\trSymplectic{}结构在此不起任何作用。
\begin{remark}\label{thm:original-definition-remark}
[131] 中所使用的\trFibrationWithSingularities{}定义，把我们的第三条条件（$TE^h$ 的\trHorizontal{}性）换成下面这个更严格的版本：
\begin{itemize}
\item \emph{在 $\partial^h E$ 附近的平凡性。}设 $M=E_*$ 是某根\trFiber。那么存在 $B\times\partial M$ 的一个\trOpen{}邻域 $V'\subset B\times M$、$\partial^h E$ 的一个\trOpen{}邻域 $V\subset E$，以及一个逐\trFiber\trDiffeomorphism{}$V'\to V$——它在 $V'\cap(\{*\}\times M)$ 上是恒同——使得 $\theta_E$ 的\trPullback{}为 $\theta_B+\theta_M$，而 $I_E$ 的\trPullback{}为乘积\trAlmostComplex{}结构 $I_B\times I_M$。
\end{itemize}

关于 $\theta_E$ 的这个条件意味着\trConnection{}$TE^h$ 在 $\partial^h E$ 附近是平凡的，并且（在底空间的\trRegular{}部分之上）给出了\trStructureGroup{}到 $\mathrm{Aut}(M,\partial M)$ 的一个\trCanonical{}约化。这条性质使得这类\trFibration{}处理起来稍微容易一些；但另一方面，也存在未能满足它的自然构造。
\end{remark}
